\hspace{3mm}

\centerline{\bf \Huge Игры и стратегии}

\centerline{\bf \large Базовый уровень}

\begin{flushright}
        \scriptsize{}
\end{flushright}

\problem{0.1} 
Двое по очереди ломают шоколадку $6 \times 8$ За ход разрешается сделать
прямолинейный разлом любого из кусков вдоль углубления. Проигрывает тот, кто не сможет сделать ход.

\problem{0.2} 
Двое по очереди кладут пятаки на круглый стол, причем так, чтобы они не накладывались друг на друга. Проигрывает тот, кто не может сделать ход.

\problem{0.3}
Двое по очереди ставят слонов в клетки шахматной доски так, чтобы слоны не били друг друга. (Цвет слонов значения не имеет). Проигрывает тот, кто не может сделать ход.

\problem{1}
На доске написаны $10$ единиц и $10$ двоек. За ход разрешается стереть две любые цифры и, если они были одинаковыми, написать двойку, а если разными --- единицу. Если последняя оставшаяся на доске цифра --- единица, то выиграл первый игрок, если двойка --- то второй.


\problem{2}
Имеется три кучки камней: в первой --- $10$, во второй --- $15$, в третьей --- $20$. За ход разрешается разбить любую кучку на две меньшие; проигрывает тот, кто не сможет сделать ход.


\problem{3}
Двое по очереди ставят коней в клетки шахматной доски так, чтобы кони не били друг друга. Проигрывает тот, кто не может сделать ход.


\problem{4}
Двое по очереди ставят королей в клетки доски 9 × 9 так, чтобы короли не били друг друга. Проигрывает тот, кто не может сделать ход.

\problem{5}
На столе лежат 40 бабуфохов. Герман и Эрдем могут взять из этой кучки либо 1 либо 2 бабуфоха. Проигрывает тот, кто не сможет ничего взять. Кто выиграет, если Герман ходит первым?

\newpage

\hspace{3mm}

\centerline{\bf \Huge Игры и стратегии}

\centerline{\bf \large Продвинутый уровень}

\begin{flushright}
        \scriptsize{}
\end{flushright}

\problem{0.1} Ладья стоит на поле $a1$. За ход разрешается сдвинуть её на любое число клеток вправо или на любое число клеток вверх. Выигрывает тот, кто поставит ладью на $h8.$ У кого есть выигрышная стратегия?

\problem{1} 
В кастрюле лежит 2027 беригов. Герман и Ердем по очереди съедают из кастрюли от 1 до 17 беригов. Когда кастрюля опустошается, игроки подсчитывают съеденное количество беригов. Если эти числа взаимно просты – выигрывает Ердем, в противном случае выигрывает Герман. Ердем ходит первым.

\problem{2}
Король стоит на поле $a1$ шахматной доски. За ход разрешается сдвинуть его на одну клетку вправо, или на одну клетку вверх, или на одну клетку вправо-вверх. Выигрывает тот, кто поставит короля на клетку $h8$. Кто выигрывает при правильной игре?

\problem{3} 
Имеются две кучки конфет: в одной --- $20$, в другой --- $21$. За ход нужно съесть одну из кучек, а вторую разделить на две не обязательно равных кучки. Проигрывает тот, кто не может сделать ход.

\problem{4}
На столе лежат 2002 карточки с числами $1, 2, 3,... , 2002$. Двое играющих берут по одной карточке по очереди. После того, как будут взяты все карточки, выигравшим считается тот, у кого больше последняя цифра суммы чисел на взятых карточках. Кто из играющих может всегда выигрывать, как бы ни играл противник, и как он должен при этом играть?

\problem{5} 
 Герман и Анджур Аркадьевич играют в следующую игру. Герман выписывает подряд цифры 0 или 1 в произвольном порядке, пока не получится 2023-значное число. Каждый раз, когда Герман выставляет цифры, Анджур Аркадьевич меняет местами две уже написанные цифры(первый ход он пропускает, а остальные обязательно ходит!) Если итоговое число получится палиндромом(то есть будет читаться с обеих сторон одинаково), то Анджур Аркадьевич побеждает, если нет, то побеждает Герман.

 \problem{6}
 Ферзь стоит на поле $c1$. За ход его можно передвинуть на любое число полей вправо, вверх или по диагонали "вправо-вверх". Выигрывает тот, кто поставит ферзя на поле $h8$.
